\subsection{(Co)lax idempotent monads}
\label{subsec:prelim-monads}

\gutentag{00B8}%
Monads on $ (\infty, 2) $-categories are defined using
Lurie's formalism of homotopy coherent algebra, with the help of
the functor
$ \iota _{1} \colon \twoCat \to \Cat $
discarding non-invertible $ 2 $-cells.

\begin{defi}
\label{def:monad}
\gutentag{00B9}%
A \emph{monad} on $\D$ is an associative algebra $T$ in the monoidal
$\infty$-category
$\iota_{1}\func(\D,\D)$
\cite[Def.~4.1.1.6]{ha}, with unit
$\eta\colon \mathrm{id}_{\D} \Rightarrow T$
and multiplication
$\mu\colon TT \Rightarrow T$.
Its $\infty$-category of \emph{modules} is
$\modcat{T}{\D} := \mathrm{LMod}_{T}(\iota_{1}\D)$
\cite[Def.~4.2.1.13]{ha}, with forgetful functor
$U\colon \modcat{T}{\D} \to \iota_{1}\D$.
A module consists of an object $x$, an action $a\colon Tx \to x$ and higher
coherence data; abusing notation, we write it as $(x,a)$,
and call the $1$-morphisms of $\modcat{T}{\D}$ \emph{module maps}.
For lax idempotent $T$, we will refine $\modcat{T}{\D}$ to a locally full
sub-$(\infty,2)$-category of $\D$ in \cref{thm:kz-module-maps}.
\todo{Upgrade $\modcat{T}{\D}$ to an $(\infty,2)$-category for general $T$,
e.g.\ as the Eilenberg--Moore object of \cite{riehlverity2016adjunctions}, a
weighted limit in $\twoCat$.}
A \emph{comonad} on $\D$ is a monad on $\D\opcat$.
\end{defi}

\gutentag{00BA}%
Only the underlying $\infty$-category of $\func(\D,\D)$ enters: the unit and
associativity laws hold up to coherent invertible modifications,
so in $2$-category theory $T$ would be called a \emph{pseudomonad}
\cite{marmolejo1997doctrines}. We will import results on pseudomonads whose
proofs construct $2$-morphisms from the unit and associativity $2$-morphisms of
$T$, of modules and of module maps, and compare composites of them. Such proofs
apply verbatim, and the $2$-morphisms satisfying the triangle identities that
they produce extend to adjunctions by \cref{thm:adj-lifting}(2).
\todo{2-categorical input: check, for each imported result, that the proof
works verbatim.}
The $\infty$-category $\modcat{T}{\D}$ only depends on the image of $T$ in
$\func(\iota_{1}\D,\iota_{1}\D)$,
a monad on $\iota_{1}\D$ in the sense of
\cite[Def.~4.7.0.1]{ha}. The functor $U$ is conservative
\cite[Cor.~4.2.3.2]{ha}, it has the left adjoint
$F\colon x \mapsto (Tx,\mu_{x})$
\cite[Cor.~4.2.4.8]{ha}, and every $U$-split simplicial object of
$\modcat{T}{\D}$ has a colimit, which $U$ preserves \cite[Thm.~4.7.3.5]{ha}.

\begin{defi}
\label{def:kz}
\gutentag{00BB}%
A monad $T$ on $\D$ is \emph{lax idempotent} if there is an adjunction in
$\func(\D,\D)$ with left adjoint $T\eta$, right adjoint $\mu$ and unit the unit
law
$\mathrm{id}_{T} \simeq \mu \circ T\eta$.
Via the equivalence
$\iota_{1}\func(\D^{\mathrm{co}},\D^{\mathrm{co}}) \simeq
\iota_{1}\func(\D,\D)$,
$T$ is a monad on $\D^{\mathrm{co}}$, with
$\modcat{T}{\D^{\mathrm{co}}} = \modcat{T}{\D}$,
and $T$ is \emph{colax idempotent} if it is lax idempotent as a monad on
$\D^{\mathrm{co}}$.
A comonad on $\D$ is lax, resp.\ colax, idempotent if it is so as a monad on
$\D\opcat$.
\end{defi}

\gutentag{00BC}%
Here $(-)^{\mathrm{co}}$ is an automorphism of $\twoCat$ preserving products, so
$\func(\D^{\mathrm{co}},\D^{\mathrm{co}}) \simeq \func(\D,\D)^{\mathrm{co}}$
compatibly with composition and evaluation, and $\iota_{1}$ does not see
$(-)^{\mathrm{co}}$.
Having the adjunction of \cref{def:kz} is a condition on $T$, not a structure:
by \cref{thm:adj-lifting}(3), the space of adjunctions
$T\eta \dashv \mu$
with unit the unit law is empty or contractible.
Lax idempotent monads are also called
Kock--Z\"oberlein monads \cite{kock1995monads}; we follow the terminology of
\cite{kellylack1997propertylike}.

\begin{lem}
\label{lem:kz-basic}
\gutentag{00BD}%
Let $T$ be a monad on $\D$.
\begin{enumerate}
\item $T$ is lax idempotent if and only if there is an adjunction with left
adjoint $\mu$, right adjoint $\eta T$ and counit the unit law
$\mu \circ \eta T \simeq \mathrm{id}_{T}$.
A lax idempotent monad thus has adjunctions
$T\eta \dashv \mu \dashv \eta T$,
with the unit laws as the unit of the first and the counit of the second.
Evaluating at $x \in \D$ gives
$T\eta_{x} \dashv \mu_{x} \dashv \eta_{Tx}$.
\item $T$ is colax idempotent if and only if there is an adjunction with left
adjoint $\mu$, right adjoint $T\eta$ and counit the unit law, if and only if
there is an adjunction with left adjoint $\eta T$, right adjoint $\mu$ and unit
the unit law.
A colax idempotent monad thus has adjunctions
$\eta T \dashv \mu \dashv T\eta$,
with the unit laws as the unit of the first and the counit of the second.
Evaluating at $x \in \D$ gives
$\eta_{Tx} \dashv \mu_{x} \dashv T\eta_{x}$.
\item Via
$\iota_{1}\func(\D\opcat,\D\opcat) \simeq (\iota_{1}\func(\D,\D))\opcat$,
a comonad on $\D$ is a coassociative coalgebra $G$ in
$\iota_{1}\func(\D,\D)$,
with counit
$\varepsilon\colon G \Rightarrow \mathrm{id}_{\D}$
and comultiplication
$\delta\colon G \Rightarrow GG$.
It is lax, resp.\ colax, idempotent if and only if there are adjunctions
\begin{equation*}
\varepsilon G \dashv \delta \dashv G\varepsilon \ , \qquad \text{resp.} \qquad
G\varepsilon \dashv \delta \dashv \varepsilon G \ ,
\end{equation*}
with the counit laws as the counit of the first and the unit of the second
adjunction; either adjunction suffices.
\item Let $\E$ be an $(\infty,2)$-category. Composition makes
$T_{\natural} := T \circ -$
a monad on $\func(\E,\D)$ and
$T^{\natural} := - \circ T$
a monad on $\func(\D,\E)$. If $T$ is lax idempotent, then $T_{\natural}$ is lax
idempotent and $T^{\natural}$ is colax idempotent.
\end{enumerate}
\end{lem}

\begin{proof}
\gutentag{00BE}%
(1) For pseudomonads in $2$-categories this is
\cite[Thm.~11.1]{marmolejo1997doctrines} (see
\cite[Thm.~6.2]{kellylack1997propertylike} for $2$-monads); we give the
argument. Composition makes $\func(\D,\D)$ an associative algebra in $\twoCat$
\cite[\S 4.7.1]{ha}, so whiskering with $T$ on either side is a functor of
$(\infty,2)$-categories, and $\eta$ is natural:
$\eta S' \circ \alpha \simeq T\alpha \circ \eta S$
for
$\alpha\colon S \Rightarrow S'$.
Let
$T\eta \dashv \mu$
with unit the unit law and counit $\varepsilon$. Let
$\delta\colon T\eta \Rightarrow \eta T$
and
$\nu\colon \mathrm{id}_{TT} \Rightarrow \eta T \circ \mu$
be the composites
\begin{equation*}
T\eta \simeq T\eta \circ \mu \circ \eta T \xRightarrow{\ \varepsilon\eta T\ }
\eta T \ , \qquad \mathrm{id}_{TT} \simeq T\mu \circ T\eta T
\xRightarrow{\ T\mu \circ \delta T\ } T\mu \circ \eta TT \simeq \eta T \circ
\mu \ ,
\end{equation*}
using the unit law
$\mu \circ \eta T \simeq \mathrm{id}_{T}$
and the naturality of $\eta$ at $\mu$.
By the triangle identity
$\mu\varepsilon \simeq \mathrm{id}_{\mu}$,
$\mu\delta$ is the composite of the unit laws
$\mu \circ T\eta \simeq \mathrm{id}_{T} \simeq \mu \circ \eta T$,
and by the triangle identity
$\varepsilon T\eta \simeq \mathrm{id}_{T\eta}$,
$\delta\eta$ is the naturality equivalence
$T\eta \circ \eta \simeq \eta T \circ \eta$.
As
$\mu \circ T\mu \simeq \mu \circ \mu T$,
$\mu\nu$ is $(\mu\delta)T$ and
$\nu\eta T$ is
$T\mu \circ (\delta\eta)T$,
up to the coherences of the unit and associativity laws; so $\nu$ and the unit
law
$\mu \circ \eta T \simeq \mathrm{id}_{T}$
satisfy the triangle identities, and
$\mu \dashv \eta T$
with this counit by \cref{thm:adj-lifting}(2). The converse is the same
argument for $T$ as an associative algebra in
$\func(\D,\D)^{\mathrm{co}}$
with
the reversed composition, which exchanges $T\eta$ and $\eta T$, and left and
right adjoints.
Evaluation at $x$ is a functor
$\func(\D,\D) \to \D$
taking $T\eta$, $\mu$ and $\eta T$ to $T\eta_{x}$, $\mu_{x}$ and $\eta_{Tx}$.

(2) An adjunction $l \dashv r$ in an $(\infty,2)$-category $\A$ is an adjunction
$r \dashv l$ in $\A^{\mathrm{co}}$, with unit and counit exchanged
(\cref{thm:adj-lifting}(2)). As
$\func(\D^{\mathrm{co}},\D^{\mathrm{co}}) \simeq \func(\D,\D)^{\mathrm{co}}$,
(2) is \cref{def:kz} and (1) for $T$ on $\D^{\mathrm{co}}$.

(3) Likewise
$\func(\D\opcat,\D\opcat) \simeq \func(\D,\D)\opcat$
compatibly with composition,
$\iota_{1}(\A\opcat) = (\iota_{1}\A)\opcat$,
and
$l \dashv r$ in $\A$ is an adjunction
$r\opcat \dashv l\opcat$
in $\A\opcat$ with the same unit and counit
\cite[Rem.~4.3]{haugseng2021laxmonads}.
This takes the adjunctions of (1) and (2) for a monad on $\D\opcat$ to the
ones stated.

(4) By \cite[\S 4.7.1]{ha}, composition makes
$S \mapsto S_{\natural}$
a monoidal functor
$\func(\D,\D) \to \func(\func(\E,\D),\func(\E,\D))$,
and
$S \mapsto S^{\natural}$
a monoidal functor on $\func(\D,\D)$ with the reversed composition.
Both carry $T$ to monads, and as functors of $(\infty,2)$-categories they
preserve adjunctions. As precomposition reverses the order of composition,
they take
$T\eta \dashv \mu$
to
$T_{\natural}\eta_{\natural} \dashv \mu_{\natural}$
and
$\eta^{\natural}T^{\natural} \dashv \mu^{\natural}$,
with the unit laws as units. So $T_{\natural}$ is lax idempotent, and
$T^{\natural}$ is colax idempotent by (2).
\end{proof}

\begin{prop}
\label{prop:kz-action}
\gutentag{00BF}%
Let $T$ be a lax idempotent monad on $\D$ and $(x,a)$ a $T$-module. Then there
is an adjunction
$a \dashv \eta_{x}$
with counit the unit law
$a\eta_{x} \simeq \mathrm{id}_{x}$;
in particular $\eta_{x}$ is reflective (\cref{def:reflection}). Dually, if $T$
is colax idempotent, then there is an adjunction
$\eta_{x} \dashv a$
with unit the unit law, and $\eta_{x}$ is coreflective.
\end{prop}

\begin{proof}
\gutentag{00C0}%
This is \cite[Lemma~10.4]{marmolejo1997doctrines}
(\cite[Thm.~6.2]{kellylack1997propertylike} for $2$-monads).
\todo{2-categorical input: check that the argument works verbatim.} The unit is
\begin{equation*}
\mathrm{id}_{Tx} \simeq Ta \circ T\eta_{x} \xRightarrow{\ Ta\,\delta_{x}\ } Ta
\circ \eta_{Tx} \simeq \eta_{x}a \ ,
\end{equation*}
where
$\delta\colon T\eta \Rightarrow \eta T\mu T\eta \simeq \eta T$
is given by the unit of
$\mu \dashv \eta T$
(\cref{lem:kz-basic}(1)). The colax case follows by \cref{def:kz}.
\end{proof}

\begin{thm}[being a module is a property]
\label{thm:kz-modules}
\gutentag{00C1}%
Let $T$ be a lax idempotent monad on $\D$. An object $x$ of $\D$ admits a
$T$-module structure if and only if $\eta_{x}$ is reflective, and the space of
$T$-module structures on $x$, the fibre of $U^{\simeq}$ over $x$, is then
contractible; the action is a left adjoint of $\eta_{x}$. Dually, for $T$ colax
idempotent, with $\eta_{x}$ coreflective and the action a right adjoint of
$\eta_{x}$.
\end{thm}

\gutentag{00C2}%
We will prove \cref{thm:kz-modules} at the end of \cref{subsec:mates-mates},
together with the recognition of module maps by mates
(\cref{thm:kz-module-maps}): two module structures on $x$ will be compared by
lifting $\mathrm{id}_{x}$ to a module map.
